\documentclass[11pt,a4paper]{article}
\usepackage[utf8]{inputenc}
\usepackage[T1]{fontenc}
\usepackage[spanish]{babel}
\usepackage{hyperref}
\usepackage{geometry}
\geometry{margin=2.5cm}
\usepackage{parskip}

\title{333 la visión es el otro que es muy diferente y ayuda}
\author{Alejandro Valenzuela Rivera\\\small Universidad Autónoma de Nuevo León}
\date{\today}
\begin{document}
\maketitle

\begin{quote}
\textbf{Sesgo declarado.} Este trabajo se sostiene desde el sesgo declarado de su autor, y no finge neutralidad. Lo que afirma se puede tumbar, y ese es el punto: el falsificador está escrito, no escondido.
\end{quote}

\section*{Resumen}
Una visión es el que es muy diferente, aprendió algo que nosotros no, y nos ayuda por algo. La alineación se explica; la visión se muestra.

\section{Por qué se sostiene}
Lo que se explica se puede copiar; lo que se muestra no: hay que haberlo visto.

\section{De dónde viene}
La autoridad pidió 'visiones' en vez de 'alineaciones', y dijo por qué.

\section{El caso}
Visión: el androide del Mind Stone que se vuelve fiable porque ve lo que los demás no.

\section{Método}
\textquestiondown{}Se muestra, o solo se explica?

\section{Límites, y qué lo tumba}
Una visión que no trae nada que nosotros no supiéramos.

\section{Lo que queda abierto}
\textquestiondown{}La ayuda viene de alguien que sabe algo que nosotros no?

\section*{Cómo se leyó esta tesis}
Esto se escribió en el repositorio man-made-empathy, del que este candado es una pieza: candados/06 del caso/333 la visión es el otro que es muy diferente y ayuda.md. La máquina que lo compuso lo hizo desde este mismo archivo, y no agregó nada que no estuviera escrito. Si algo de esto es falso, se dice en el mismo sitio, y se corrige con fecha.

\end{document}
